\begin{pastexam}[label={px:ntu-111-bime-ds-04}]{臺大生機所 111 學年度（C）第 4 題}
（證明）

(10\%) Prove that, if a connected graph of $N$ nodes has the property that removing any edge
disconnects the graph, then the graph has $N-1$ edges and no cycles.
\end{pastexam}

\begin{pastsoln}
圖指無向圖。 拿掉以後會讓圖不連通的邊叫做\textbf{橋}（bridge，資結書第~23.5~節）；題目的條件是
「$G$ 連通，而且每一條邊都是橋」。 分兩步：先證沒有環，再由「連通、無環」推出邊數是 $N - 1$。

\textbf{第 1 步：$G$ 沒有環。} 用反證法。 假設 $G$ 有一個環 $v_1, v_2, \ldots, v_k, v_1$（$k \ge 3$），取環上
的一條邊 $e = \{v_1, v_2\}$，把它拿掉得到 $G - e$。 任取兩個頂點 $x, y$，在 $G$ 裡有一條 $x$ 到 $y$ 的路徑
（$G$ 連通）。 如果這條路徑沒有經過 $e$，它在 $G - e$ 裡還在；如果經過 $e$，把 $e$ 這一步換成沿著環的
另一半 $v_1, v_k, v_{k-1}, \ldots, v_2$ 走（這幾條邊都不是 $e$），得到 $G - e$ 裡一條 $x$ 到 $y$ 的路線。 所以
$G - e$ 仍然連通，$e$ 不是橋，和條件矛盾。 因此 $G$ 沒有環。

\textbf{第 2 步：連通、無環、$N$ 個頂點的圖有 $N - 1$ 條邊。} 對 $N$ 歸納。

$N = 1$：只有一個頂點，沒有環就不能有自環，$0 = N - 1$ 條邊。

$N \ge 2$：先證有一個度數為 $1$ 的頂點（葉）。 取圖中一條最長的簡單路徑 $P = u_0, u_1, \ldots, u_t$
（圖連通且 $N \ge 2$，所以 $t \ge 1$）。 $u_0$ 的鄰居除了 $u_1$ 不可能有別的：鄰居 $w$ 若在 $P$ 上（$w = u_i$，
$i \ge 2$），$u_0, u_1, \ldots, u_i, u_0$ 是一個環；若不在 $P$ 上，$w, u_0, \ldots, u_t$ 是更長的簡單路徑。 兩者都
不可能，所以 $\deg(u_0) = 1$。

把 $u_0$ 和它唯一的邊拿掉，得到 $G'$：$N - 1$ 個頂點；仍然無環（只少了東西）；仍然連通（任兩個留下來的
頂點之間的簡單路徑不會經過度數 $1$ 的 $u_0$，因為經過一個頂點要用到它的兩條邊）。 由歸納假設 $G'$ 有
$N - 2$ 條邊，加回 $u_0$ 的那條邊，$G$ 有 $N - 1$ 條邊。 \hfill$\blacksquare$

所以 $G$ 是一棵樹：連通、無環、$N - 1$ 條邊，這正是資結書定理~21.3 裡「樹」的幾種等價說法。 反過來也
成立：樹的每一條邊都是橋（拿掉 $\{u, v\}$ 後若 $u$、$v$ 仍相連，那條路加上 $\{u, v\}$ 就是環）。
\end{pastsoln}
